\section{Conclusion}
\label{sec:conclusion}

We began by observing that benchmark evaluation produces aggregate scores without explaining why models fail, blocking targeted correction. When an LLM fails a TruthfulQA question, existing tools report the failure but provide no diagnostic signal distinguishing entity-substitution errors from reasoning failures. This gap prevents matched correction strategies --- practitioners apply RAG or COT uniformly, missing opportunities to target root causes.

Our work demonstrates that attention entropy over entity spans provides a measurable diagnostic signal (p < 0.001, Cohen's d = -1.13), enabling automated failure-type classification with 86.7\% accuracy. Entity-substitution errors exhibit zero entropy in 60\% of cases, revealing they are precision errors (model looks at wrong place deterministically) rather than recall errors (absence of focused attention). This mechanistic understanding grounds matched correction routing: RAG retrieves the correct entity for attention re-direction, targeting the entity-level misdirection root cause.

Our main contributions are:

\begin{itemize}
\item \textbf{Methodological:} Attention entropy-based failure classification at scale (73 processed samples vs 10-20 manual examples in prior interpretability work), demonstrating that automated diagnostic routing from benchmark to interpretability analysis is tractable.

\item \textbf{Empirical:} Robust attention pattern signature (p = $7.5 \times 10^{-7}$, large effect size d = -1.13) with zero-entropy entity-errors revealing precision error mechanism, validated through statistical testing and threshold-based classification.

\item \textbf{Framework (structural feasibility only):} Mock validation with synthetic RAG/COT success rates demonstrates matched routing structure (+24pp over mismatched), but \textbf{correction effectiveness is unproven} --- real-world Wikipedia API + GPT-judge validation is HIGH-priority future work (FW2).
\end{itemize}

We acknowledge principled limitations bounding claims to proof-of-concept: (1) GPT-2-only attention patterns --- multi-model replication pending (FW1 HIGH), (2) \textbf{synthetic correction results} --- real-world RAG effectiveness unvalidated (FW2 HIGH, critical for credibility), (3) manual gold labels limiting scale (FW4), (4) single failure type --- reasoning errors untested (FW3), (5) 27\% span alignment loss (FW5). We validate diagnostic classification robustly; correction routing awaits real-world deployment.

\subsection{Future Directions}

Immediate next steps test pattern generalization and real-world correction effectiveness. \textbf{FW1} (HIGH priority) replicates attention pattern detection with Llama-2-7B, GPT-3.5, and GPT-4 to determine if the entropy difference generalizes across architectures. \textbf{FW2} (HIGH priority) implements full RAG pipeline (Wikipedia API + GPT-3.5 generation) and real COT baseline, evaluating correction effectiveness on 50 entity-errors using GPT-judge. This experiment is critical for paper credibility --- without real-world validation, we can claim diagnostic classification but not correction routing effectiveness.

Longer-term extensions explore multi-failure-type diagnosis, automated labeling, and cross-benchmark generalization. \textbf{FW3} (MEDIUM priority) tests whether reasoning errors and knowledge gaps exhibit distinct attention signatures, enabling multi-class classification. \textbf{FW4} (MEDIUM priority) trains supervised classifiers on gold-labeled samples to automate failure categorization at scale, removing the manual annotation bottleneck.

As LLM deployment scales, the gap between ``knowing a model failed'' and ``understanding why it failed'' becomes a reliability bottleneck. Attention-based failure routing bridges that gap, scaling interpretability from manual analysis to automated diagnosis. Our framework demonstrates that benchmarks, interpretability, and correction can form an integrated system --- transforming evaluation from passive measurement to active improvement guidance.
